\FloatBarrier
\section{Discussion and Conclusion}
\label{sec:current-headroom}

The forward, backward, and training results point to three hardware
takeaways.  We state each takeaway first and then give the measurement that
supports it; the rejected kernel variants remain in the appendices.

\begin{table}[htbp]
\centering
\caption{Hardware takeaways and the evidence that supports them.}
\label{tab:hardware-takeaways}
\small
\begin{tabularx}{\textwidth}{p{.28\textwidth}X}
\toprule
Takeaway & Supporting evidence \\
\midrule
Tensor-memory ownership limits overlap & Two FP32 score banks and two FP32
output accumulators use all $4\times128=512$ tensor-memory columns; reducing
shared-memory use did not increase CTA residency or reduce latency. \\
Readiness stalls leave matrix throughput unused & A matched diagnostic kept
the same 98,304 tensor instructions but reached only 18.8\% tensor-pipe
activity with the real probability path. \\
Probability arithmetic is no longer the main gap & A fixed-P diagnostic that
removes nearly all probability construction improved latency by only 5.23\%. \\
\bottomrule
\end{tabularx}
\end{table}

\subsection{Takeaway 1: TMEM capacity and ownership bound the overlap window}

The two-query pipeline in Figure~\ref{fig:quarter-pipeline} already overlaps
QK, softmax, and PV across query stages.  The limitation is more specific:
under the current D128 layout, the next QK cannot fully overlap the still-owned
probability and tail-PV phase.  Figure~\ref{fig:tmem-layout} shows why.  Two
128-column score banks and two 128-column FP32 output accumulators occupy all
512 tensor-memory (TMEM) columns.  A retired part of a score bank is reused for
the quantized probability and scale pages.  Until PV consumes that overlay,
the next QK targeting that score bank and query stage has no legal score
destination; the other query stage still provides partial overlap.

This is a storage-ownership dependency, not simply a shortage of barriers or
shared memory.  In a matched control, reducing dynamic shared memory from
209,920 to 163,840 bytes did not improve latency because the 512-column TMEM
allocation still allowed only one CTA per streaming multiprocessor.  More
shared memory, or more raw on-chip bytes that cannot be allocated as another
score bank, would not change this schedule.

\begin{figure}[htbp]
\centering
\begin{tikzpicture}[font=\small,
  bank/.style={draw=inkgray,minimum width=2.35cm,minimum height=.72cm,
    align=center},
  extra/.style={draw=scoreblue,dashed,minimum width=2.35cm,
    minimum height=.72cm,align=center}]
  \node[anchor=east,font=\bfseries] at (-.25,1.05) {Current};
  \node[bank,fill=scoreblue!22] (s0) at (1.0,1.05) {score 0 / P 0};
  \node[bank,fill=scoreblue!34,right=0pt of s0] (s1) {score 1 / P 1};
  \node[bank,fill=outputgreen!22,right=0pt of s1] (o0) {output 0};
  \node[bank,fill=outputgreen!34,right=0pt of o0] (o1) {output 1};
  \node[below=3pt of s1.south east,anchor=north,text=inkgray]
    {all 512 columns owned};

  \node[anchor=east,font=\bfseries] at (-.25,-.55) {Candidate contract};
  \node[bank,fill=scoreblue!22] (t0) at (1.0,-.55) {score 0 / P 0};
  \node[bank,fill=scoreblue!34,right=0pt of t0] (t1) {score 1 / P 1};
  \node[bank,fill=outputgreen!22,right=0pt of t1] (u0) {output 0};
  \node[bank,fill=outputgreen!34,right=0pt of u0] (u1) {output 1};
  \node[extra,fill=scoreblue!8,right=4pt of u1] (next) {next score};
  \node[below=4pt of t1.south,font=\scriptsize,text=scoreblue]
    {PV still consuming};
  \node[above=4pt of next.north,font=\scriptsize,text=scoreblue] (qkwrite)
    {next QK writes};
  \draw[-{Latex[length=2mm]},thick,scoreblue]
    (qkwrite.south) -- (next.north);
\end{tikzpicture}
\caption{Conceptual hardware implication.  The useful addition is another
allocatable score destination, together with legal issue and ownership
semantics, so the next QK can proceed while PV still consumes the current
overlay.  The lower row is not a measured kernel or a claim that raw capacity
alone is sufficient.}
\label{fig:extra-score-bank}
\end{figure}

\subsection{Takeaway 2: low tensor activity is a readiness symptom}

The scaled-FP4 PV instruction consumes a K64 inner dimension, so two adjacent
32-column probability fragments must be ready before useful PV work can
start.  Publishing the first fragment alone does not advance the matrix
issuer.  The score-to-probability path therefore creates paired rendezvous:
Q0 and Q1 enable the first PV half, while Q2 and Q3 enable the tail.  Attempting
to issue the next QK between the first and tail PV operations violates the
current sequence of two K64 tensor-core accumulations.

Table~\ref{tab:historical-readiness-profile} shows a matched historical
profiling checkpoint.  The real and fixed-probability variants execute the
same number of tensor instructions.  Removing probability work shortens wall
time and raises the fraction of cycles attributed to the tensor pipe, but it
does not create more matrix work.  Source sampling attributes 1,715 of 2,669
not-issued samples to dependency stalls categorized by NVIDIA's profiler as
long-scoreboard waits, chiefly final statistics, score readiness, and output
publication.  Tensor-core underutilization is therefore a symptom of operands
not being ready, not a lack of nominal FP4 throughput.

\begin{table}[htbp]
\centering
\caption{Matched historical profile with identical tensor work.  This
diagnostic predates the final Direct-P binary and is used only to identify the
stall mechanism.}
\label{tab:historical-readiness-profile}
\small
\begin{tabular}{lrrr}
\toprule
Probability path & Dynamic instructions & Tensor instructions & Tensor active \\
\midrule
Real probability construction & 54.6 million & 98,304 & 18.8\% \\
Fixed probability & 11.9 million & 98,304 & 26.6\% \\
\bottomrule
\end{tabular}
\end{table}

The final Direct-P path has already removed most exposed probability cost.
At B1/S4096/H24/D128, the retained kernel measures 0.092448 ms and repeats at
0.092512 ms.  Intentionally incorrect ceilings bound what remains:

\begin{table}[htbp]
\centering
\caption{Final-kernel ceilings relative to the 0.092448-ms valid record.  These
rows deliberately remove required work and do not compute valid attention.}
\label{tab:speed-of-light}
\small
\begin{tabular}{lrr}
\toprule
Diagnostic & Time (ms) & Gap from 0.092448 ms \\
\midrule
Simplified score packing & 0.091168 & 1.280\us\ (1.38\%) \\
Keep row maximum, pack raw scores & 0.090112 & 2.336\us\ (2.53\%) \\
Use a fixed probability tile & 0.087616 & 4.832\us\ (5.23\%) \\
\bottomrule
\end{tabular}
\end{table}

Even the fixed-P diagnostic that removes nearly all probability construction
saves only 5.23\%.  The remaining gap to the four-times-BF16 matrix ratio
includes score loads, K64 publication granularity, scale delivery, online
correction, output accumulation, and the epilogue.  This is also why fewer
arithmetic instructions do not always help: independent arithmetic can occupy
issue slots while the kernel is waiting for a dependency.

\subsection{Takeaway 3: faster arithmetic does not scale the whole kernel}

\begin{table}[htbp]
\centering
\caption{Hardware properties relevant to the measured kernels.  TMEM capacity
does not increase from GB200 to the tested B300.}
\label{tab:blackwell-generations}
\small
\begin{tabular}{lcc}
\toprule
Property & GB200 / SM100 & B300 / SM103 \\
\midrule
Visible SMs in this study & 152 & 148 \\
Maximum reported clock & 2062 MHz & 2032 MHz \\
TMEM per SM & 256 KB & 256 KB \\
Key exponential rate & 16 ops/clock/SM & 32 ops/clock/SM \\
Dense NVFP4 GPU class & 1.0$\times$ & 1.5$\times$ \\
Fused TMEM load and reduction & no & yes \\
\bottomrule
\end{tabular}
\end{table}

Blackwell Ultra increases dense NVFP4 throughput, doubles important
special-function rates, and adds a fused TMEM load/reduction instruction
\cite{nvidia2025blackwellultra,nvidia2026ptx,zadouri2026flashattention4}.
Those improvements do not enlarge the score/output allocation, so the whole
attention loop does not become 1.5$\times$ faster.  The measurements are
consistent with faster QK and PV making score reduction, scale publication,
and correction more visible.

Launch shape matters as well.  A persistent grid wider than the 148 visible
B300 SMs creates a mostly empty second wave; recompiling for the device removes
that tail.  With the corrected geometry, S8192/H64 and the wave-aligned
S9472/H64 cases sustain 3116 and 3159 TFLOP/s.  At S32768/H24, however, B300
reaches 2945 TFLOP/s versus 2998 on GB200.  Per-SM-clock normalization favors
B300, but that normalization is an inference rather than an observed
GPU-throughput win.

Noncausal D64 measurements and unsafe-format controls are reported in
Appendix~\ref{app:nvnv-d64}.  Their different tile and ownership regime should
not be inferred from the D128 hardware profile above.

\subsection{Implications for hardware and kernels}

The measured dependencies motivate these candidate changes:
\begin{description}
  \item[Another allocatable score bank.] Allow the next QK to write while PV
  still owns the current probability overlay.  This must preserve two-query
  K/V reuse and provide legal issue semantics; a naive ping-pong prototype
  doubled the query jobs and was 52.8\% slower.
  \item[K32 scaled-FP4 PV.] Consume each 32-column probability completion
  immediately instead of waiting for a K64 pair.
  \item[Scales outside TMEM.] Remove instruction-facing scale pages from the
  overlay by reading compact scales from shared memory or registers.  Existing
  scale-footprint probes show that this alone neither creates a full
  128-column score destination nor demonstrates a speedup.
  \item[Wider tiles with a compatible lifecycle.] B300's wider matrix
  instructions have high synthetic ceilings, but only if probability
  production, K/V reuse, and score storage can be scheduled together.
\end{description}

These are hypotheses motivated by measured dependencies, not measured wins.
The resulting design implication is that, once Direct-P shortens probability
construction, another polynomial approximation is likely less valuable than
a larger usable overlap window.
